%% kapa: exam=1 round=28 year=2017 period=2 form=A subject=accounting subject_ko=회계학 seq=27 number=67
%% source: Q-Net 감정평가사 기출문제 (cst00309) artlSeq=5199603 file=extracted/1008/5199603/2017년 제28회 감정평가사 1차시험 문제지/2교시 A형.pdf
%% conversion: kordoc 4.24.51 -> tools/gen1.py (auto); review=proofread-sonnet; flags=none
\begin{question}{67}
(주)감평은 20x1년 1월 1일에 건물을 ￦5,000,000에 취득(내용연수 10년, 잔존가치 ￦0, 정액법 감가상각)하였다. 20x1년 말 및 20x2년 말 기준 원가모형을 적용하는 건물의 순공정가치는 각각 ￦3,600,000과 ￦3,900,000이고, 사용가치는 각각 ￦3,000,000과 ￦4,300,000이다. (주)감평은 건물의 회수가능액과 장부금액의 차이가 중요하고 손상징후가 있는 것으로 판단하여 손상차손(손상차손환입)을 인식하였다. 관련 설명으로 옳지 않은 것은?
\begin{choices}
  \item 20x2년도에 감가상각비로 ￦400,000을 인식한다.
  \item 20x1년 말 재무상태표에 표시되는 건물 장부금액은 ￦3,600,000이다.
  \item 20x2년 말 재무상태표에 표시되는 건물 장부금액은 ￦4,000,000이다.
  \item 20x1년도에 손상차손으로 ￦900,000을 인식한다.
  \item 20x2년도에 손상차손환입으로 ￦1,100,000을 인식한다.
\end{choices}
\end{question}
\answer{5}
